\section{Satélite}
Os sistemas de satélite funcionam, na sua maior parte, como um caminho de reflexão e amplificação para o sinal que está sendo transmitido. A transmissão e recepção se dão em visada direta, contornando problemas de curvatura da Terra para pontos que estão muito distantes entre si e evitando problemas de desvanecimento e sombreamento. Porém as grandes distâncias envolvidas impõem uma alta taxa de atenuação do sinal e consequentemente um requerimento de potência de transmissão do sinal bem mais alto. Na maioria das aplicações é necessário que o Satélite seja capaz de regenerar o sinal recebido antes de retransmitir para o receptor. Cada sistema de satélite requer uma certa alocação de banda para seu funcionamento, a depender da aplicação(TV, Voz, Internet, Radio ou Coleta de Dados).  

\section{Celular, 3G e 4G}
A comunicação celular 3G consiste de uma série de padrões de comunicação que atingem certos critérios técnicos IMT-2000, dentre estes estão UMTS, W-CDMA, TD-SCDMA, HSPA+, CDMA2000 e EDGE. A comunicação usando os padrões 3G e 4G depende da alocação geográfica de torres de celular, que formam "clusters" e distribuições hexagonais, a partir das quais é possível fazer a alocação de canais para cada usuário. 

\section{Wifi, Família IEEE 802.11}
